\section{Conclusion}

We developed and evaluated an end-to-end fraud detection pipeline for the ULB
credit card dataset, treating the problem as a business decision problem rather
than a pure classification exercise.

\textbf{Methodologically}, we established a leakage-controlled workflow: a
single persisted stratified split, resampling confined to training folds
inside an \texttt{imblearn} pipeline, and scaler fitting inside the pipeline.
Across 27 configurations spanning five model families and five imbalance
techniques, Random Forest with SMOTE at a 1:100 ratio emerged as the strongest
configuration, attaining F1 = 0.8605 and Average Precision = 0.8239 on the
held-out test set.

\textbf{Empirically}, we produced several findings that qualify common
practice. Expanding the hyper-parameter search degraded test performance while
improving cross-validated scores --- consistently across both model families,
and consistent with overfitting of the search itself on a small minority class,
though a single seed and a small CV gap do not establish this as a proof. Tuning the operating threshold on the test set
inflated recall by 4.9 percentage points, which we quantified rather than
assumed away. Logistic Regression's failure at the 95\% precision constraint
traced to model capacity rather than to the class ratio. Random Undersampling,
one of the most commonly taught techniques, was the most expensive approach and
the only one that degraded every model family it was paired with. Finally, we
showed that Random Forest's default impurity importance mis-ranks the single
most informative feature, an error that would have gone undetected without
permutation importance and SHAP.

\textbf{Practically}, we converted predictive output into an operating policy,
and this turned out to matter more than the model choice. A single threshold
cannot simultaneously satisfy the 95\% precision floor and the
100-review-per-day capacity that the brief imposes: at threshold 0.52 the
workload is 193 reviews per day (1.9 analysts), while the tightest threshold
that fits within one analyst's capacity collapses recall to 40.9\%. A
\textbf{two-tier decision architecture} --- automatic blocking above 0.810,
human review between 0.175 and 0.810, approval below --- resolves the
dilemma: it attains \textbf{81.0\% recall} at a total cost of USD 154.65 per
85{,}118 transactions while requiring only \textbf{70 analyst reviews per
day}, one analyst's worth of work. Because auto-blocking consumes no analyst
time, the two tiers can be tuned independently, and the result is an
operational improvement rather than a re-labelling: the comparable recall
reached by the cheapest single threshold near that operating point (0.1767,
recall 0.8028) would need 232 reviews per day, more than twice the capacity. The model requires
0.045 ms per transaction and under USD 2.30 per year in compute even at ten
million transactions per day, so the constraint on deployment is
infrastructure rather than modelling.

\textbf{Finally}, we compared the model against four rule-based baselines. It
caught more than twice as many frauds than the best rule while flagging 99.0\%
fewer transactions, at one twenty-third of the cost. The contribution of
machine learning here is not a marginally better metric; it is a fundamentally
cheaper way of protecting both the bank and its customers.

The honest headline result is that, at a 95\% precision constraint, our model
catches roughly three quarters of all fraud, at a cost of USD 170 per 85{,}118
transactions, and that the remaining quarter consists almost entirely of very
small transactions that the available features cannot distinguish from
ordinary purchases. Closing that gap requires better data, not a better
algorithm.
